% Preamble for the dotnet-rs proof-DSL feasibility study.
% Compiled with XeLaTeX (TeX Live 2026). Fonts loaded by filename via kpathsea
% because the TeX Live OTFs are not registered with system fontconfig.

\usepackage{fontspec}
\usepackage{amsmath}
\usepackage{amsthm}
\usepackage{unicode-math}

\setmainfont{LibertinusSerif}[
  Extension      = .otf,
  UprightFont    = *-Regular,
  BoldFont       = *-Bold,
  ItalicFont     = *-Italic,
  BoldItalicFont = *-BoldItalic,
]
\setsansfont{LibertinusSans}[
  Extension   = .otf,
  UprightFont = *-Regular,
  BoldFont    = *-Bold,
  ItalicFont  = *-Italic,
]
\setmonofont{LibertinusMono}[
  Extension    = .otf,
  UprightFont  = *-Regular,
  BoldFont     = *-Regular,
  BoldFeatures = {FakeBold=1.4},
  ItalicFont   = *-Regular,
  ItalicFeatures = {FakeSlant=0.2},
  Scale        = MatchLowercase,
]
\setmathfont{LibertinusMath-Regular.otf}

\usepackage[margin=1.1in]{geometry}
\usepackage{microtype}
\setlength{\emergencystretch}{2.5em}
\usepackage{booktabs}
\usepackage{array}
\usepackage{longtable}
\usepackage{enumitem}
\usepackage{xcolor}
\usepackage{graphicx}
\usepackage{mathpartir}
\usepackage{listings}
\usepackage[most]{tcolorbox}
\usepackage{caption}
\usepackage{csquotes}
\usepackage[
  backend=bibtex,
  style=numeric,
  maxbibnames=6,
]{biblatex}
\usepackage[hidelinks,unicode]{hyperref}
\usepackage{cleveref}

% ---------------------------------------------------------------------------
% Colors
% ---------------------------------------------------------------------------
\definecolor{rustkw}{RGB}{146, 46, 14}
\definecolor{dslkw}{RGB}{20, 60, 120}
\definecolor{commentgray}{RGB}{100, 100, 100}
\definecolor{stringgreen}{RGB}{30, 100, 50}
\definecolor{boxframe}{RGB}{60, 60, 90}
\definecolor{boxbg}{RGB}{247, 247, 250}
\definecolor{verdictbg}{RGB}{240, 246, 240}
\definecolor{verdictframe}{RGB}{50, 100, 60}
\definecolor{riskbg}{RGB}{250, 244, 240}
\definecolor{riskframe}{RGB}{140, 80, 40}

% ---------------------------------------------------------------------------
% Listings: Rust, CIL, the proposed DSL, and Lean
% ---------------------------------------------------------------------------
\lstdefinelanguage{Rust}{
  keywords={as,break,const,continue,crate,dyn,else,enum,extern,false,fn,for,
    if,impl,in,let,loop,match,mod,move,mut,pub,ref,return,self,Self,static,
    struct,super,trait,true,type,unsafe,use,where,while,async,await},
  keywordstyle=\color{rustkw}\bfseries,
  ndkeywords={u8,u16,u32,u64,usize,i8,i16,i32,i64,isize,f32,f64,bool,char,
    str,String,Vec,Option,Some,None,Result,Ok,Err,Box,Arc,Rc},
  ndkeywordstyle=\color{dslkw},
  comment=[l]{//},
  morecomment=[s]{/*}{*/},
  commentstyle=\color{commentgray}\itshape,
  string=[b]",
  stringstyle=\color{stringgreen},
  sensitive=true,
}
\lstdefinelanguage{ProofDSL}{
  keywords={invariant,requires,ensures,proof,qed,by,from,have,show,obtain,
    assume,witness,axiom,lemma,forall,exists,let,in,at,region,token,resource,
    open,apply,frame,rooted,live,aligned,valid,init,disjoint,stw,holds,
    justified,spec,rust,obligation},
  keywordstyle=\color{dslkw}\bfseries,
  comment=[l]{--},
  commentstyle=\color{commentgray}\itshape,
  string=[b]",
  stringstyle=\color{stringgreen},
  sensitive=true,
  mathescape=true,
}
\lstdefinelanguage{Lean}{
  keywords={theorem,lemma,def,structure,inductive,where,by,exact,intro,apply,
    fun,match,with,do,let,have,show,sorry,namespace,end,open,import,variable,
    instance,class,extends,deriving,axiom,opaque,partial,termination_by},
  keywordstyle=\color{dslkw}\bfseries,
  comment=[l]{--},
  morecomment=[s]{/-}{-/},
  commentstyle=\color{commentgray}\itshape,
  sensitive=true,
  mathescape=true,
}
\lstdefinelanguage{CIL}{
  keywords={ldarg,ldloc,stloc,ldfld,stfld,ldflda,ldind,stind,ldobj,stobj,
    cpblk,initblk,cpobj,initobj,newobj,newarr,call,callvirt,ret,box,unbox,
    volatile,unaligned,constrained,readonly,ldelema,ldelem,stelem,localloc},
  keywordstyle=\color{rustkw},
  comment=[l]{//},
  commentstyle=\color{commentgray}\itshape,
  sensitive=true,
}
\lstset{
  basicstyle=\ttfamily\small,
  columns=fullflexible,
  keepspaces=true,
  breaklines=true,
  breakatwhitespace=true,
  showstringspaces=false,
  frame=single,
  framerule=0.3pt,
  rulecolor=\color{boxframe!40},
  xleftmargin=1em,
  framexleftmargin=0.6em,
  aboveskip=0.9\baselineskip,
  belowskip=0.9\baselineskip,
  captionpos=b,
}

% ---------------------------------------------------------------------------
% Boxes
% ---------------------------------------------------------------------------
\newtcolorbox{verdictbox}[1][]{
  enhanced, breakable,
  colback=verdictbg, colframe=verdictframe,
  boxrule=0.5pt, arc=1.5pt,
  left=8pt, right=8pt, top=6pt, bottom=6pt,
  fonttitle=\bfseries\sffamily,
  title={#1},
}
\newtcolorbox{riskbox}[1][]{
  enhanced, breakable,
  colback=riskbg, colframe=riskframe,
  boxrule=0.5pt, arc=1.5pt,
  left=8pt, right=8pt, top=6pt, bottom=6pt,
  fonttitle=\bfseries\sffamily,
  title={#1},
}
\newtcolorbox{obligationbox}[1][]{
  enhanced, breakable,
  colback=boxbg, colframe=boxframe,
  boxrule=0.5pt, arc=1.5pt,
  left=8pt, right=8pt, top=6pt, bottom=6pt,
  fonttitle=\bfseries\sffamily,
  title={#1},
}

% ---------------------------------------------------------------------------
% Theorem environments
% ---------------------------------------------------------------------------
\theoremstyle{definition}
\newtheorem{definition}{Definition}[section]
\newtheorem{example}{Example}[section]
\theoremstyle{plain}
\newtheorem{proposition}{Proposition}[section]
\newtheorem{obligation}{Obligation}[section]
\theoremstyle{remark}
\newtheorem{remark}{Remark}[section]

% ---------------------------------------------------------------------------
% Semantic macros
% ---------------------------------------------------------------------------
\newcommand{\code}[1]{\texttt{#1}}
\newcommand{\crate}[1]{\textsf{#1}}
\newcommand{\dotnetrs}{\code{dotnet-rs}}
\newcommand{\gcarena}{\code{gc-arena}}
\newcommand{\ecma}[1]{\textsection{}#1}
\newcommand{\safetyc}{\code{//\,SAFETY:}}

% Judgments and semantic objects
\providecommand{\llbracket}{[\mkern-3.5mu[}
\providecommand{\rrbracket}{]\mkern-3.5mu]}
\newcommand{\sem}[1]{\llbracket #1 \rrbracket}
\newcommand{\ok}{\ \mathsf{ok}}
\newcommand{\rooted}{\mathsf{rooted}}
\newcommand{\live}{\mathsf{live}}
\newcommand{\alignedp}{\mathsf{aligned}}
\newcommand{\valid}{\mathsf{valid}}
\newcommand{\init}{\mathsf{init}}
\newcommand{\stw}{\mathsf{stw}}
\newcommand{\hoare}[3]{\{\,#1\,\}\;#2\;\{\,#3\,\}}
\renewcommand{\wp}[2]{\mathsf{wp}\;#1\;\{\,#2\,\}}
\newcommand{\later}{\mathop{\triangleright}}
\newcommand{\sep}{\mathrel{*}}
\newcommand{\wand}{\mathrel{-\!\!*}}
\newcommand{\pointsto}{\mapsto}
\newcommand{\vsafe}{\vdash_{\mathsf{safe}}}
\newcommand{\TCB}{\mathrm{TCB}}
